% Zusammenhängender Prosabeweis; Aussagen und Kennungen bleiben gemeinsam.
\Needspace{6\baselineskip}
\begin{proof}[Beweis]
Seien \(F\colon A\inj B\) und \(G\colon B\inj A\) gegeben.
Wir konstruieren eine Bijektion, indem wir zwei Einschränkungen dieser
Funktionen passend zusammensetzen. Die Buchstaben \(S,C,D\) und
\(\mathcal K\) bezeichnen dabei nur Hilfsmengen dieses Beweises.

\medskip\noindent
\textbf{Die kleinste abgeschlossene Teilmenge.}
Setze
\[
  \begin{aligned}
    S&:=A\setminus G[B],\\
    \mathcal K&:=\{X\in\powerset(A)\mid
      S\subseteq X\land G[F[X]]\subseteq X\}.
  \end{aligned}
\]
Die Familie \(\mathcal K\) ist durch Aussonderung aus der Potenzmenge
von \(A\) eine Menge. Sie ist nicht leer: Da \(F[A]\subseteq B\)
und \(G[B]\subseteq A\) gelten, erfüllt \(A\) beide Bedingungen und
gehört selbst zu \(\mathcal K\).

Wir definieren \(C\) als den Durchschnitt dieser Familie:
\[
  C=\{x\in A\mid \forall X\in\mathcal K\;x\in X\}
   =\bigcap\mathcal K.
\]
Der Durchschnitt ist hier über eine nichtleere Mengenfamilie gebildet.
Da \(A\in\mathcal K\) gilt, ist er eine Teilmenge von \(A\). Jedes Mitglied von
\(\mathcal K\) enthält \(S\); folglich gilt \(S\subseteq C\).
Sei nun \(X\in\mathcal K\) beliebig. Aus \(C\subseteq X\)
folgt durch Monotonie der Bildbildung
\(G[F[C]]\subseteq G[F[X]]\). Weil \(X\) zu \(\mathcal K\)
gehört, gilt nach deren Definition \(G[F[X]]\subseteq X\).
Zusammen ergibt dies
\[
  G[F[C]]\subseteq G[F[X]]\subseteq X.
\]
Jeder Punkt aus \(G[F[C]]\) liegt somit in allen Mitgliedern von
\(\mathcal K\), also in \(C\). Damit ist \(C\) selbst ein Mitglied
von \(\mathcal K\). Zugleich ist \(C\) eine Teilmenge jedes anderen Mitglieds
und ist daher die kleinste Teilmenge von \(A\), welche die beiden
geforderten Bedingungen erfüllt.

\medskip\noindent
\textbf{Die Fixpunktgleichung.}
Aus der Minimalität ergibt sich die stärkere Aussage
\[
  C=S\cup G[F[C]].
\]
Setze dazu \(D:=S\cup G[F[C]]\). Die gerade bewiesenen Inklusionen
liefern \(D\subseteq C\), also insbesondere \(D\subseteq A\).
Wiederum wegen der Monotonie der Bildbildung gilt
\[
  G[F[D]]\subseteq G[F[C]]\subseteq D.
\]
Auch \(S\subseteq D\) gilt nach Definition. Somit gehört \(D\)
zu \(\mathcal K\), und die Minimalität von \(C\) liefert
\(C\subseteq D\). Zusammen mit \(D\subseteq C\) ist die
Fixpunktgleichung bewiesen. Sie beschreibt \(C\) als Vereinigung
von \(S\) und dem Bild von \(C\) unter \(G\circ F\).

\medskip\noindent
\textbf{Das Bild des zweiten Zweiges.}
Die Fixpunktgleichung besagt
\(C=(A\setminus G[B])\cup G[F[C]]\). Beim Übergang zum
Komplement in \(A\) entsteht nach der De-Morganschen Regel zunächst
ein Schnitt:
\[
  \begin{aligned}
    A\setminus C
      &=\bigl(A\setminus(A\setminus G[B])\bigr)
        \cap\bigl(A\setminus G[F[C]]\bigr)\\
      &=G[B]\cap\bigl(A\setminus G[F[C]]\bigr)\\
      &=G[B]\setminus G[F[C]].
  \end{aligned}
\]
Die zweite Zeile benutzt \(G[B]\subseteq A\): Ein Punkt von \(A\),
der nicht in \(A\setminus G[B]\) liegt, liegt gerade in \(G[B]\).
In der letzten Zeile kann die Bedingung, zu \(A\) zu gehören,
entfallen, weil jeder Punkt von \(G[B]\) sie bereits erfüllt.
Es bleiben also genau die Punkte von \(G[B]\), die nicht zu
\(G[F[C]]\) gehören.

Auch \(G[B]\setminus(A\cap G[F[C]])\) ist dieselbe Menge:
Aus \(F[C]\subseteq B\) folgt
\(G[F[C]]\subseteq G[B]\subseteq A\), also
\(A\cap G[F[C]]=G[F[C]]\). Der zusätzliche Schnitt mit \(A\)
ändert die Menge somit nicht.

Mit der Injektivität von \(G\) lässt sich die rechte Seite auch
als Bild einer Differenz schreiben:
\[
  G[B]\setminus G[F[C]]=G[B\setminus F[C]].
\]
Dies lässt sich unmittelbar an den Elementen prüfen. Ist
\(b\in B\setminus F[C]\), so gehört \(G(b)\) zu \(G[B]\).
Läge \(G(b)\) auch in \(G[F[C]]\), gäbe es ein \(c\in F[C]\)
mit \(G(b)=G(c)\). Die Injektivität von \(G\) ergäbe \(b=c\),
im Widerspruch zu \(b\notin F[C]\). Umgekehrt besitzt jeder Punkt
aus \(G[B]\setminus G[F[C]]\) ein Urbild \(b\in B\).
Dieses Urbild kann nicht zu \(F[C]\) gehören und liegt deshalb in
\(B\setminus F[C]\). Insgesamt haben wir damit gezeigt:
\[
  G[B\setminus F[C]]=A\setminus C.
\]

\medskip\noindent
\textbf{Die beiden Bijektionen.}
Eine injektive Funktion ist auf jeder Teilmenge ihres Definitionsbereichs
bijektiv auf ihr Bild. Wegen \(C\subseteq A\) ist daher
\[
  F\restriction_C\colon C\bij F[C]
\]
eine Bijektion. Entsprechend ist die Einschränkung von \(G\) auf
\(B\setminus F[C]\) bijektiv auf dessen Bild. Die gerade bewiesene
Bildgleichung identifiziert dieses Bild mit \(A\setminus C\):
\[
  G\restriction_{B\setminus F[C]}\colon
  B\setminus F[C]\bij A\setminus C.
\]
Erst diese Bijektion kehren wir um und erhalten
\[
  J:=\bigl(G\restriction_{B\setminus F[C]}\bigr)^{-1}
  \colon A\setminus C\bij B\setminus F[C].
\]

\medskip\noindent
\textbf{Zusammensetzen.}
Definiere die Funktion \(H\) auf \(A\) durch
\[
  H(a):=
  \begin{cases}
    F(a),&a\in C,\\
    J(a),&a\in A\setminus C.
  \end{cases}
\]
Die beiden Definitionsbereiche sind disjunkt und vereinigen sich zu
\(A\); jeder Zweig nimmt seine Werte in \(B\) an. Damit ist \(H\)
eine wohldefinierte Funktion von \(A\) nach \(B\), also gerade die
Fallfunktion \(\CaseFun{C}{A\setminus C}{F\restriction_C}{J}\).

Innerhalb jedes Zweiges ist \(H\) injektiv. Auch zwischen den Zweigen
können keine gleichen Werte auftreten: Der erste Zweig hat das Bild
\(F[C]\), der zweite das dazu disjunkte Bild \(B\setminus F[C]\).
Folglich ist \(H\) auf ganz \(A\) injektiv. Jeder Punkt aus \(B\)
gehört zu einem dieser beiden Bildbereiche und wird vom zugehörigen
bijektiven Zweig getroffen. Somit ist \(H\) auch surjektiv und daher
die gesuchte Bijektion \(H\colon A\bij B\).
\end{proof}

\CSBHeading{Einordnung}
Der vorstehende Beweis benötigt weder eine Folge natürlicher Zahlen noch eine
Rekursionskonstruktion oder das Auswahlaxiom. Auch leere Mengen sind
eingeschlossen: Die Familie \(\mathcal K\) bleibt nichtleer, weil
sie \(A\) enthält. Die Hilfsmenge \(C\) dient der Konstruktion;
für spätere Anwendungen genügt das Ergebnis, dass zwei Injektionen
in entgegengesetzten Richtungen die Existenz einer Bijektion sichern.
